\begin{proof}[Proof of \cref{obj:lem:template-conditioning}]
\begin{enumerate}
\item Fix \(Q\in\mathcal Q_h\). Since \(h=2^{-j}>0\), feasibility gives
\[
0<h^{-2\beta}\le N_{Q\ell}
\qquad\text{for every }\ell\in\{0,\ldots,m\}^d.
\]
Thus every microcell count is positive, and its reciprocal is well defined. We establish the conditioning bounds using these positive counts.
% lean: CausalSmith.Stat.WeakOverlap.gramHat_minEig_ge

\item Define the coefficient and node index sets, and the polynomial associated with a coefficient vector, by
\[
\mathcal I_m=\{\alpha\in\mathbb N_0^d:|\alpha|\le m\},
\qquad
\mathcal T_m=\{0,\ldots,m\}^d,
\qquad
r_b(z)=\sum_{\alpha\in\mathcal I_m}b_\alpha z^\alpha
\quad(z\in\mathbb R^d).
\]
Expanding the definition of \(G_0\) yields
\[
b^\top G_0b=\frac1J\sum_{\ell\in\mathcal T_m}r_b(v_\ell)^2\ge0.
\]

The nodes in \cref{obj:synth_1} form the Cartesian product of the one-dimensional set
\[
\mathcal S_{\mathrm{node}}
=\left\{\frac{a+1}{m+2}:a=0,\ldots,m\right\},
\qquad
|\mathcal S_{\mathrm{node}}|=m+1.
\]
Each coordinate degree of \(r_b\) is at most \(m\). Consequently, vanishing at every tensor node implies that \(r_b\) is the zero polynomial. Indeed, fixing all but one coordinate at node values gives a univariate polynomial with \(m+1\) distinct roots and degree at most \(m\), so it vanishes identically. Repeating this argument coordinate by coordinate extends the vanishing to all of \(\mathbb R^d\). Hence
\[
\bigl[r_b(v_\ell)=0\text{ for every }\ell\in\mathcal T_m\bigr]
\Longrightarrow r_b\equiv0
\Longrightarrow b_\alpha=0\text{ for every }\alpha\in\mathcal I_m.
\]
For \(d=0\), both index sets have one element and \(r_b\) is constant, so the same implication follows directly. A unit coefficient vector has a nonzero coordinate; therefore at least one summand is strictly positive, and
\[
\|b\|_2=1\quad\Longrightarrow\quad b^\top G_0b>0.
\]

To obtain strict positivity of the infimum, define
\[
\mathcal S_{\mathrm{coef}}
=\left\{b\in\mathbb R^{\mathcal I_m}:
\sum_{\alpha\in\mathcal I_m}b_\alpha^2=1\right\},
\qquad
(b_{\mathrm{const}})_\alpha=
\begin{cases}
1,&\alpha=0,\\
0,&\alpha\ne0.
\end{cases}
\]
The zero multi-index belongs to \(\mathcal I_m\) for every \(d,m\), and
\(\|b_{\mathrm{const}}\|_2=1\). Thus \(\mathcal S_{\mathrm{coef}}\) is nonempty.
It is closed, and each of its coordinates satisfies
\[
|b_\alpha|^2\le\sum_{\alpha'\in\mathcal I_m}b_{\alpha'}^2=1,
\]
so it is bounded and hence compact. The continuous quadratic form attains its minimum there: for some \(b_{\min}\in\mathcal S_{\mathrm{coef}}\),
\[
\lambda_0
=\min_{b\in\mathcal S_{\mathrm{coef}}}b^\top G_0b
=b_{\min}^\top G_0b_{\min}>0.
\]
% lean: CausalSmith.Stat.WeakOverlap.templateLambda_pos

\item Fix an arbitrary \(b\in\mathbb R^{\mathcal I_m}\). Write the supremum norm, including its zero-dimensional convention, as
\[
\|z\|_\infty
=\max\bigl(\{0\}\cup\{|z_a|:1\le a\le d\}\bigr).
\]
The template Lipschitz constant described in \cref{obj:synth_2} is
\[
L_U=
\sup_{\substack{z,z'\in[0,1]^d\\z\ne z'}}
\frac{\|U(z)U(z)^\top-U(z')U(z')^\top\|_{\mathrm{op}}}
{\|z-z'\|_\infty},
\]
with value zero when the set of quotients is empty. For \(d\ge1\),
\cref{obj:synth_2} supplies its finiteness; for \(d=0\), the cube is a singleton and the stated convention applies. Every quotient is nonnegative, so \(L_U\ge0\).

By \cref{obj:def:template}, the width and microcubes satisfy
\[
\eta=\min\left\{\frac1{2(m+2)},\frac{\lambda_0}{1+L_U}\right\}>0,
\qquad
V_\ell=v_\ell+[-\eta/2,\eta/2]^d.
\]
In particular,
\[
\frac{\eta}{2}\le\frac1{4(m+2)}\le\frac1{m+2},
\qquad
\frac1{m+2}\le(v_\ell)_a
\le1-\frac1{m+2}.
\]
The coordinate inequalities are vacuous when \(d=0\). They imply, for every \(z\in V_\ell\),
\[
v_\ell,z\in[0,1]^d,
\qquad
\|z-v_\ell\|_\infty\le\frac{\eta}{2}.
\]

Define the perturbation matrix and the coefficient-dependent error by
\[
M_\ell(z)
=U(v_\ell)U(v_\ell)^\top-U(z)U(z)^\top,
\qquad
\varepsilon_b=L_U\frac{\eta}{2}\|b\|_2^2.
\]
If \(z=v_\ell\), the desired squared-evaluation inequality follows immediately from \(\varepsilon_b\ge0\). If \(z\ne v_\ell\), the definition of \(L_U\) and the preceding distance bound give
\[
\frac{\|M_\ell(z)\|_{\mathrm{op}}}{\|v_\ell-z\|_\infty}
\le L_U,
\qquad
\|M_\ell(z)\|_{\mathrm{op}}
\le L_U\|v_\ell-z\|_\infty
\le L_U\frac{\eta}{2}.
\]
For \(b=0\), the quadratic-form bound below is immediate. For \(b\ne0\), define
\[
b_{\mathrm{unit}}=\frac{b}{\|b\|_2},
\qquad
\|b_{\mathrm{unit}}\|_2=1.
\]
The operator-norm definition and homogeneity give
\[
\|M_\ell(z)b\|_2
=\|b\|_2\|M_\ell(z)b_{\mathrm{unit}}\|_2
\le\|M_\ell(z)\|_{\mathrm{op}}\|b\|_2.
\]
Cauchy--Schwarz therefore yields
\[
\begin{aligned}
r_b(v_\ell)^2-r_b(z)^2
&=b^\top M_\ell(z)b\\
&\le\|b\|_2\|M_\ell(z)b\|_2\\
&\le\|M_\ell(z)\|_{\mathrm{op}}\|b\|_2^2
\le\varepsilon_b.
\end{aligned}
\]
Combining the two cases for \(z\), we obtain
\[
r_b(v_\ell)^2-\varepsilon_b\le r_b(z)^2
\qquad(z\in V_\ell).
\]
% lean: CausalSmith.Stat.WeakOverlap.templatePointwise_square_perturbation

\item Define the normalized observations and the treated index sets by
\[
z_{Qi}=\frac{X_i-a_Q}{h}\quad(1\le i\le n),
\qquad
\mathcal I_{Q\ell}
=\{i\in\{1,\ldots,n\}:A_i=1,\ X_i\in E_{Q\ell}\}.
\]
Then \(|\mathcal I_{Q\ell}|=N_{Q\ell}>0\). Since \(h>0\) and
\(E_{Q\ell}=a_Q+hV_\ell\), every \(i\in\mathcal I_{Q\ell}\) satisfies
\(z_{Qi}\in V_\ell\). Averaging the pointwise inequality within each cell gives
\[
\begin{aligned}
r_b(v_\ell)^2-\varepsilon_b
&=\frac1{N_{Q\ell}}\sum_{i\in\mathcal I_{Q\ell}}
\bigl(r_b(v_\ell)^2-\varepsilon_b\bigr)\\
&\le\frac1{N_{Q\ell}}\sum_{i\in\mathcal I_{Q\ell}}r_b(z_{Qi})^2.
\end{aligned}
\]
Expansion of the empirical Gram matrix, followed by equal averaging across the \(J>0\) cells, now gives
\[
\begin{aligned}
b^\top\widehat G_Qb
&=\frac1J\sum_{\ell\in\mathcal T_m}
\frac1{N_{Q\ell}}\sum_{i\in\mathcal I_{Q\ell}}r_b(z_{Qi})^2\\
&\ge\frac1J\sum_{\ell\in\mathcal T_m}
\bigl(r_b(v_\ell)^2-\varepsilon_b\bigr)\\
&=b^\top G_0b-\varepsilon_b.
\end{aligned}
\]
The last equality uses \(|\mathcal T_m|=J\), so the equally averaged error remains \(\varepsilon_b\).
% lean: CausalSmith.Stat.WeakOverlap.templateGram_quadratic_sub_error_le_gramHat

\item For \(b=0\), the template lower bound is an equality. For \(b\ne0\), the unit vector defined above and the defining infimum give
\[
\lambda_0\le b_{\mathrm{unit}}^\top G_0b_{\mathrm{unit}},
\qquad
\lambda_0\|b\|_2^2
\le\|b\|_2^2b_{\mathrm{unit}}^\top G_0b_{\mathrm{unit}}
=b^\top G_0b.
\]
The width rule also implies
\[
\eta\le\frac{\lambda_0}{1+L_U},
\qquad
L_U\eta\le\eta(1+L_U)\le\lambda_0.
\]
Consequently,
\[
\varepsilon_b
=L_U\frac{\eta}{2}\|b\|_2^2
\le\frac{\lambda_0}{2}\|b\|_2^2.
\]
Combining these inequalities with the equal-cell transfer proves
\[
b^\top\widehat G_Qb
\ge b^\top G_0b-\varepsilon_b
\ge\lambda_0\|b\|_2^2-\frac{\lambda_0}{2}\|b\|_2^2
=\frac{\lambda_0}{2}\|b\|_2^2.
\]
% lean: CausalSmith.Stat.WeakOverlap.gramHat_minEig_ge_of_pos

\item If \(\det(\widehat G_Q)=0\), there is a nonzero coefficient vector \(b_{\ker}\) with
\[
\widehat G_Qb_{\ker}=0,
\qquad
\|b_{\ker}\|_2>0.
\]
Applying the quadratic lower bound gives the contradiction
\[
0<\frac{\lambda_0}{2}\|b_{\ker}\|_2^2
\le b_{\ker}^\top\widehat G_Qb_{\ker}=0.
\]
Thus \(\det(\widehat G_Q)\ne0\).
% lean: CausalSmith.Stat.WeakOverlap.det_ne_zero_of_quadratic_lower

\item For an arbitrary input vector \(b_{\mathrm{in}}\in\mathbb R^{\mathcal I_m}\), define
\[
b_{\mathrm{inv}}=\widehat G_Q^{-1}b_{\mathrm{in}}.
\]
Invertibility gives \(\widehat G_Qb_{\mathrm{inv}}=b_{\mathrm{in}}\). The quadratic lower bound and Cauchy--Schwarz imply
\[
\frac{\lambda_0}{2}\|b_{\mathrm{inv}}\|_2^2
\le b_{\mathrm{inv}}^\top\widehat G_Qb_{\mathrm{inv}}
=b_{\mathrm{inv}}^\top b_{\mathrm{in}}
\le\|b_{\mathrm{inv}}\|_2\|b_{\mathrm{in}}\|_2.
\]
If \(\|b_{\mathrm{inv}}\|_2=0\), the following bound is immediate; otherwise, divide by its positive norm and then by \(\lambda_0/2>0\). In both cases,
\[
\|\widehat G_Q^{-1}b_{\mathrm{in}}\|_2
=\|b_{\mathrm{inv}}\|_2
\le\frac2{\lambda_0}\|b_{\mathrm{in}}\|_2.
\]
The coefficient unit sphere is nonempty, as witnessed by \(b_{\mathrm{const}}\). Taking the supremum over its vectors therefore gives
\[
\|\widehat G_Q^{-1}\|_{\mathrm{op}}
=\sup_{\|b_{\mathrm{in}}\|_2=1}
\|\widehat G_Q^{-1}b_{\mathrm{in}}\|_2
\le\frac2{\lambda_0}.
\]
Since \(Q\) was arbitrary, all the asserted bounds hold throughout \(\mathcal Q_h\).
% lean: CausalSmith.Stat.WeakOverlap.finiteL2OpNorm_inv_le_of_quadratic_lower
\end{enumerate}
\end{proof}
